%chapters/3-deep-computations/

\chapter{Deep Computations}\label{ch:deep_computations}

The contents of this chapter are based on my contributions to and perspectives on \cite{Duenezetal2026}, the second paper in a series started by the authors of \cite{alva2024approximability} that connects model theory, structural Ramsey theory, and the topology of function spaces to the foundational study of ``deep computations.''
While the initial work in \cite{alva2024approximability} established the formal framework for ultracomputations, characterizing them as accumulation points of standard computations and proving the existence of deep equilibria, this chapter focuses on the complexity and classification of this asymptotic behaviour.

By utilizing the Todorčević trichotomy for Rosenthal compacta as a ``Rosetta stone,'' we translate abstract topological properties into a hierarchy of computational complexity classes.
We demonstrate how model-theoretic concepts, such as Shelah's classification theory and the Non-Independence Property (NIP), provide the necessary medium to distinguish between ``tame'' computations (characterized by polynomial approximability) and ``wild'' ones.
This approach provides a unified theoretical framework for learning architectures, such as Neural ODEs and Deep Equilibrium Models, by linking their asymptotic behavior to classical results in functional analysis and statistical learning theory.

% /chapters/3-deep-computations/01-intro.tex

\section{Introduction}\label{sec:deep_comp_intro}



In the space of real numbers, if $x$ is a point in the closure of some set $A\subseteq\mathbb R$, there must be a sequence of reals in $A$ that converges to $x$.
One aspect that makes studying function spaces is that this equivalence between sequential closure and topological closure need not hold.
A concrete example \cite{Duenezetal2026}
\input{chapters/3-deep-computations/02-CCS.tex}
\input{chapters/3-deep-computations/03-examples-CCS.tex}
\input{chapters/3-deep-computations/04-classifying.tex}
\input{chapters/3-deep-computations/05-randomized.tex}